\documentclass[11pt]{article}

\usepackage[margin=1.1in]{geometry}
\usepackage{amsmath}
\usepackage{booktabs}
\usepackage{tabularx}
\usepackage{array}
\usepackage{hyperref}

\title{Long-Context Retrieval in Efficient Language Models}
\author{Jeffrey Huang}
\date{May 2026}

\begin{document}

\maketitle

\section{Introduction}

Large language models increasingly depend on long contexts: source files,
research papers, transcripts, retrieved documents, and extended conversations
are now commonly presented to a model as part of a single prompt. The dominant
decoder-only Transformer architecture is well suited to such settings because
self-attention gives each token a direct mechanism for accessing earlier tokens
in the sequence \cite{vaswani2017attention}. This explicit access is one reason
Transformers perform well on retrieval-like behavior: when a later token must
recover a fact stated earlier in the prompt, attention can, in principle, route
information directly from the relevant earlier position.

The same mechanism is also the source of the long-context bottleneck. In
standard self-attention, pairwise interactions are computed between positions in
the sequence. During training this creates a quadratic dependence on context
length, while autoregressive inference requires storing and reading a growing
key-value cache. For short contexts this cost is manageable; for long contexts,
the cache becomes increasingly memory- and bandwidth-intensive. As a result, a
large body of recent work asks whether the useful behavior of attention can be
approximated, compressed, or partially replaced by more efficient sequence
mixing mechanisms.

This paper surveys several representative approaches to that problem and then
tests a narrower set of public checkpoints on synthetic long-context retrieval.
The survey is organized around a central tradeoff. Full attention preserves
explicit token-level access to the prompt, but it is expensive. Linear,
recurrent, state-space, and convolutional alternatives reduce this cost by
representing the past through structured or fixed-size state. That compression
improves efficiency, especially during inference, but it may weaken precise
retrieval when many similar facts compete for memory. Needle-in-a-haystack
tasks provide a useful probe of this tradeoff because they isolate a simple
question: can the model recover a specific key-value association from a long,
mostly irrelevant context?

\section{The Long-Context Bottleneck}

In a Transformer layer, an input sequence representation
\[
X \in \mathbb{R}^{n \times d}
\]
is projected into query, key, and value matrices,
\[
Q=XW_Q,\qquad K=XW_K,\qquad V=XW_V.
\]
Causal self-attention computes
\[
\mathrm{Attn}(X)=
\mathrm{softmax}\!\left(\frac{QK^\top+M}{\sqrt{d_k}}\right)V,
\]
where \(M\) is a causal mask. The matrix \(QK^\top\) contains all pairwise
query-key scores, so its size grows as \(n^2\) with the sequence length. This
quadratic structure is computationally expensive, but it also gives attention
its distinctive modeling strength: each token can form a content-dependent
connection to any earlier token.

Autoregressive inference changes the computational form but not the basic
length dependence. A decoder does not usually recompute the entire prefix at
each generation step. Instead, it stores previous keys and values in a
key-value cache. The current token contributes a new query, which attends over
the cached prefix. This avoids redundant computation, but the cache still grows
with context length and each new token must read from a longer memory. Long
contexts therefore pressure both computation and memory movement.

Efficient long-context architectures can be understood as different responses
to this bottleneck. Some retain attention but restrict which positions can
interact. Others replace attention with recurrent state, state-space dynamics,
linearized attention, long convolutions, or hybrid combinations of these
mechanisms. The central question for retrieval is not only whether these models
are asymptotically cheaper, but what information is lost when the full prompt is
compressed into a smaller or more structured memory.

\section{Taxonomy of Efficient Long-Context Models}

The efficient long-context literature is broad, so a chronological account can
be misleading. Rather than treating recent models as a single linear
progression, it is more useful to group them by how they represent past tokens.
Table~\ref{tab:model-landscape} summarizes the representative model families
considered in this report.

\begin{table*}[t]
\centering
\scriptsize
\setlength{\tabcolsep}{3.5pt}
\renewcommand{\arraystretch}{1.18}
\begin{tabularx}{\textwidth}{
>{\raggedright\arraybackslash}p{2.5cm}
>{\raggedright\arraybackslash}p{2.9cm}
>{\raggedright\arraybackslash}X
>{\raggedright\arraybackslash}X}
\toprule
\textbf{Model or family} &
\textbf{Memory strategy} &
\textbf{Efficiency idea} &
\textbf{Retrieval-relevant risk} \\
\midrule
Transformer &
Full causal attention &
Preserves explicit token-to-token access, but uses quadratic sequence mixing
and a growing KV cache. &
Strong retrieval in principle, but expensive at long context lengths. \\
\midrule
BigBird &
Sparse attention &
Restricts attention to local, global, and random connections to reduce the cost
of long-range attention \cite{zaheer2020bigbird}. &
Retrieval may depend on whether the sparse pattern gives a strong path to the
relevant token. \\
\midrule
RetNet &
Retention / recurrent state &
Uses a retention mechanism with parallel and recurrent forms, replacing a
growing cache with a fixed recurrent state \cite{sun2023retnet}. &
Decay-like memory may weaken old but relevant information. \\
\midrule
RWKV &
RNN-style language modeling &
Combines recurrent inference with Transformer-like training and scaling goals
\cite{peng2023rwkv}. &
Past context is summarized into recurrent state rather than explicitly stored. \\
\midrule
Hyena &
Implicit long convolution &
Uses long implicit filters and gating as an alternative to attention-based
sequence mixing \cite{poli2023hyena}. &
Long-range information is mixed through filters rather than retrieved by direct
indexing. \\
\midrule
Mamba &
Selective state-space model &
Makes state-space updates input-dependent and implements them with a
hardware-aware selective scan \cite{gu2023mamba}. &
Selective memory can retain salient inputs, but capacity is still finite. \\
\midrule
Mamba-2 &
State-space duality &
Recasts selective state-space computation into dual recurrent and structured
matrix forms for better hardware efficiency \cite{dao2024mamba2}. &
Improves the capacity-efficiency tradeoff, but still compresses history. \\
\midrule
Gated DeltaNet &
Gated delta-rule memory &
Adds a more explicit overwrite-style update to reduce interference among stored
associations \cite{yang2024gateddelta}. &
Retrieval depends on whether the learned overwrite rule preserves the needed
association under distractors. \\
\midrule
BASED &
Hybrid linear attention &
Combines global linear attention, local sliding attention, and convolutional
mixing \cite{arora2024based}. &
Compressed global memory may confuse similar key-value records. \\
\midrule
Griffin, Jamba, Nemotron-H &
Modern hybrids &
Mix recurrent, state-space, or linear components with attention layers
\cite{de2024griffin,lieber2024jamba,nvidia2025nemotron}. &
Hybrid design may recover some explicit access, but behavior depends on where
attention is retained. \\
\bottomrule
\end{tabularx}
\caption{Representative strategies for efficient long-context language
modeling. The table emphasizes how each family represents prior context and
what kind of retrieval failure that representation may induce.}
\label{tab:model-landscape}
\end{table*}

Sparse attention models such as BigBird preserve attention as the basic
operation, but reduce cost by limiting the attention pattern
\cite{zaheer2020bigbird}. This keeps some explicit token access while avoiding
the full quadratic matrix. Recurrent and state-space models instead replace
explicit history with a state updated across time. RetNet uses a retention
mechanism with recurrent and parallel forms \cite{sun2023retnet}, RWKV pursues
RNN-style language modeling at large scale \cite{peng2023rwkv}, and Mamba
introduces input-dependent selective state updates \cite{gu2023mamba}. Mamba-2
then reformulates selective state-space computation through state-space duality,
improving the practical hardware profile of this family
\cite{dao2024mamba2}.

Other approaches replace attention with different structured sequence-mixing
operations. Hyena uses implicit long convolutions and gating to model long
sequences without attention \cite{poli2023hyena}. Linear attention methods
rewrite attention-like computation so that key-value information can be
accumulated into compressed statistics; this perspective is closely related to
the view of linear Transformers as fast-weight programmers
\cite{schlag2021fastweight}. BASED builds on this direction with a hybrid
architecture combining global linear attention, local attention, and
convolution \cite{arora2024based}. More recent systems such as Griffin, Jamba,
and Nemotron-H suggest that the field is moving toward hybrid architectures
rather than pure replacements: attention remains useful, but it is mixed with
recurrent or state-space components to reduce cost
\cite{de2024griffin,lieber2024jamba,nvidia2025nemotron}.

\section{Mechanistic Themes}

Several recurring ideas cut across these model families. The first is the
distinction between explicit access and compressed state. Full attention retains
a separate key and value for every prior token. In contrast, linear attention,
state-space models, and recurrent models summarize the prefix into a fixed-size
or structured state. This shift is the main source of their efficiency, but it
also creates the possibility of interference. If many similar facts are written
into the same compressed memory, a later query may recover the wrong
association or a value that was never present in the prompt.

The second theme is forgetting. Retention and recurrent models often use
decay-like behavior to keep memory bounded and suppress stale information. Such
decay is useful when most old tokens are irrelevant, but retrieval tasks expose
the opposite case: an early token may remain essential even after many
irrelevant updates. Long-context retrieval therefore tests whether a model can
preserve a small piece of information without allowing it to fade or be
overwritten.

The third theme is selectivity. Mamba's central contribution is to make the
state update depend on the input token, allowing the model to decide which
information should be retained or suppressed \cite{gu2023mamba}. Mamba-2 keeps
this selective-memory perspective but reorganizes the computation into a form
that is more compatible with efficient matrix operations \cite{dao2024mamba2}.
Selectivity is important for retrieval because it gives the model a possible
mechanism for filtering irrelevant context. However, selectivity does not
remove the finite-state bottleneck: the model must still decide what to store
and what to sacrifice.

The fourth theme is overwrite and interference. Gated DeltaNet explicitly
targets the problem that additive memory updates can be too blunt when many
similar associations are present \cite{yang2024gateddelta}. Its gated
delta-rule update is designed to edit memory more precisely along the current
key direction. In the retrieval setting, this idea is directly relevant:
distractor records create many nearly identical write operations, and the model
must preserve the association requested by the final query without confusing it
with the others.

Finally, hybrid architectures reflect a practical compromise. BASED, Griffin,
Jamba, and Nemotron-H all combine mechanisms rather than relying on a single
sequence mixer. Local or sparse attention can preserve some exact token access,
while linear, recurrent, or state-space components provide cheaper global
context processing. This trend suggests that the long-context problem is not
simply a choice between attention and recurrence, but a question of how much
explicit access must be retained for a given task.

\section{Scope of the Empirical Study}

The survey above covers a wider landscape than can be fairly tested in a small
reproduction study. A meaningful empirical comparison requires public
checkpoints that are comparable in scale, training conditions, and
implementation quality. The original plan for this project centered on Mamba-2
and Gated DeltaNet, but that path raised substantial reproducibility concerns.
The exact Gated DeltaNet paper checkpoints were not publicly available in a
form that supported a clean comparison, and public alternatives differed in
source, configuration, and training assumptions. Dependency and kernel issues
also made it difficult to separate architecture from implementation artifacts.

For this reason, the experimental portion of the report narrows to a
HazyResearch checkpoint family released with the BASED paper. This family
includes an attention baseline, a Mamba baseline, and BASED models at two
scales. It therefore supports a more controlled comparison among full
attention, selective state-space modeling, and a hybrid linear-attention model.
The experiments should not be read as a reproduction of the BASED paper's
throughput claims. The optimized kernels required for faithful speed comparison
were not available in our setup, and BASED accuracy runs were performed with a
more conservative decoding path to avoid artifacts in the public fallback
implementation. Instead, the study is best understood as a controlled
long-context retrieval extension using public checkpoints.

\section{Experimental Expectations}

The survey suggests several expectations for synthetic retrieval. First,
single-needle retrieval should be relatively easy. When there is only one
record in the document, the model does not need to choose among competing
associations. Second, multi-key retrieval with distractors should be much more
diagnostic. In that setting, the model must bind the queried key to the correct
value among many same-format alternatives. This directly stresses interference
and overwrite behavior.

Full attention should have an advantage when exact key-value binding is
required, because it retains explicit access to prior tokens. Compressed-memory
models may still succeed if their state preserves the relevant association, but
they have a clearer failure mode: many similar records can collide in the
compressed state. For that reason, the most informative outcome is not only
accuracy but the type of error. A distractor value suggests that the model found
the record format but selected the wrong association. A number not present in
the document suggests a more general recall or generation failure. A response
with no number at all suggests that the model has left the intended answer mode
and reverted to ordinary continuation.

These are expectations rather than strict predictions. Architecture is only one
factor in empirical behavior. Checkpoint training, scale, prompt format,
tokenizer behavior, and implementation path all affect the final result. The
experiments that follow should therefore be interpreted as controlled probes of
retrieval behavior in available public models, not as definitive rankings of
all efficient long-context architectures.

\begin{thebibliography}{99}

\bibitem{vaswani2017attention}
A. Vaswani et al.
\newblock Attention is all you need.
\newblock In \emph{NeurIPS}, 2017.

\bibitem{zaheer2020bigbird}
M. Zaheer et al.
\newblock Big Bird: Transformers for longer sequences.
\newblock In \emph{NeurIPS}, 2020.

\bibitem{sun2023retnet}
Y. Sun et al.
\newblock Retentive network: A successor to Transformer for large language
models.
\newblock arXiv preprint, 2023.

\bibitem{peng2023rwkv}
B. Peng et al.
\newblock RWKV: Reinventing RNNs for the Transformer era.
\newblock arXiv preprint, 2023.

\bibitem{poli2023hyena}
M. Poli et al.
\newblock Hyena hierarchy: Towards larger convolutional language models.
\newblock In \emph{ICML}, 2023.

\bibitem{gu2023mamba}
A. Gu and T. Dao.
\newblock Mamba: Linear-time sequence modeling with selective state spaces.
\newblock arXiv preprint, 2023.

\bibitem{dao2024mamba2}
T. Dao and A. Gu.
\newblock Transformers are SSMs: Generalized models and efficient algorithms
through structured state space duality.
\newblock In \emph{ICML}, 2024.

\bibitem{yang2024gateddelta}
S. Yang et al.
\newblock Gated Delta Networks: Improving Mamba2 with delta rule.
\newblock arXiv preprint, 2024.

\bibitem{schlag2021fastweight}
I. Schlag, K. Irie, and J. Schmidhuber.
\newblock Linear Transformers are secretly fast weight programmers.
\newblock In \emph{ICML}, 2021.

\bibitem{arora2024based}
S. Arora et al.
\newblock Simple linear attention language models balance the recall-throughput
tradeoff.
\newblock arXiv preprint, 2024.

\bibitem{de2024griffin}
S. De et al.
\newblock Griffin: Mixing gated linear recurrences with local attention for
efficient language models.
\newblock arXiv preprint, 2024.

\bibitem{lieber2024jamba}
O. Lieber et al.
\newblock Jamba: A hybrid Transformer-Mamba language model.
\newblock arXiv preprint, 2024.

\bibitem{nvidia2025nemotron}
NVIDIA et al.
\newblock Nemotron-H: A family of accurate and efficient hybrid
Mamba-Transformer models.
\newblock arXiv preprint, 2025.

\end{thebibliography}

\end{document}
